% ============================================================
%  Section: RTGS — Resource-Aware Task Graph Shaping
%  IEEE Transactions style  |  English  |  third-person passive
%
%  Required preamble packages:
%    \usepackage{amsmath, amssymb, bm}
%    \usepackage{algorithm, algpseudocode}
%    \usepackage{enumitem}
%
%  Notation (consistent with system_description.tex):
%    \mathcal{G}_T = (\mathcal{V}_T, \mathcal{E}_T)  — Task DAG
%    \mathcal{G}_R = (\mathcal{U}, \mathcal{L})       — Resource graph
%    v_i  — task node  |  u_k  — UAV  |  t  — decision epoch
%    M    — number of UAVs  |  N_t — current task-graph cardinality
%    N_{\max} = 20          — padded node budget
%    \tilde{\phi}_k(t)      — battery-derated CPU of u_k at epoch t
%    B_{kl}(t)              — link bandwidth between u_k and u_l
%    \beta_k(t)             — battery fraction of u_k
% ============================================================

% ----------------------------------------------------------------
\section{RTGS: Resource-Aware Task Graph Shaping}
\label{sec:method}

% ----------------------------------------------------------------
\subsection{Framework Overview}
\label{subsec:overview}

Conventional UAV task-offloading methods treat the task DAG
$\mathcal{G}_T$ as a fixed input and optimise only the mapping
of tasks to UAVs.
This assumption is untenable in practice: heterogeneous UAV swarms
exhibit time-varying bandwidth, progressive battery depletion, and
spatial mobility, all of which continuously alter the cost of
executing and communicating between tasks.
A graph topology that is efficient at one epoch may become a
scheduling bottleneck moments later as link quality degrades or a UAV
loses computing capacity.

The proposed framework, \emph{Resource-Aware Task Graph Shaping}
(RTGS), reframes the problem by treating the task graph topology
itself as a \emph{control variable}.
At each decision epoch $t$, an RL agent observes the joint state of
the task workload and the physical resource layer, then selects a
structural transformation of $\mathcal{G}_T$ before scheduling takes
place.
The modified graph is handed to a deterministic, mobility-aware
greedy scheduler that translates the current topology into a concrete
task-to-UAV assignment.
The scheduling outcome and the evolving swarm state jointly generate
a feedback signal that closes the adaptive loop.

Three design principles distinguish RTGS from prior work.
\textbf{(i) Dynamic topology control}: the agent may expand
task parallelism via node splitting or consolidate workloads via node
merging, rather than being constrained to schedule a fixed graph.
\textbf{(ii) Resource-aware node selection}: a cross-graph attention
mechanism embeds each task node in the context of the current UAV
resource state, so that split and merge decisions are conditioned on
which physical nodes would execute the affected tasks and what links
connect them.
\textbf{(iii) Structurally valid node-level actions}: the action
space operates at the granularity of individual task nodes and
node pairs, with an action mask that enforces DAG acyclicity,
capacity feasibility, and graph-size bounds at every step.

The overall architecture is illustrated in Fig.~\ref{fig:framework}.

% ----------------------------------------------------------------
\subsection{Problem Formulation}
\label{subsec:mdp}

\subsubsection*{MDP Structure}

The RTGS control problem is formalised as a Markov Decision Process
$\langle \mathcal{S},\, \mathcal{A},\, \mathcal{R},\, \gamma \rangle$
where each component is defined to capture the dual dynamics of the
task graph and the physical UAV swarm.

\subsubsection*{State Space}

A state $s_t \in \mathcal{S}$ is a composite observation that encodes
three information sources:
\begin{equation}
  s_t = \bigl(
    \mathbf{X}^T_t,\; \mathbf{E}^T_t,\;
    \mathbf{X}^R_t,\; \mathbf{E}^R,\;
    \mathbf{f}_t
  \bigr),
  \label{eq:state}
\end{equation}
where the components are defined as follows.

\textit{Task-graph observation.}
$\mathbf{X}^T_t \in \mathbb{R}^{N_{\max} \times 2}$ is a
zero-padded node-feature matrix whose $i$-th row is the attribute
pair $\mathbf{x}_i^T = (c_i,\, s_i)$ of task $v_i$ (CPU demand and
data volume, respectively); $\mathbf{E}^T_t \in
\mathbb{Z}^{2 \times N_{\max}^2}$ is the corresponding COO-format
edge index.
Rows beyond the current graph size $N_t \leq N_{\max}$ are padded
with zeros and are excluded from all computations by structural
masking.

\textit{Resource-graph observation.}
$\mathbf{X}^R_t \in [0,1]^{M \times 3}$ is a dynamic node-feature
matrix whose $k$-th row captures the instantaneous state of UAV $u_k$:
\begin{equation}
  \mathbf{x}_k^R(t) =
  \Bigl[\,
    \tilde{\phi}_k(t)/\phi^{\max},\quad
    \bar{B}_k(t)/B_{\max},\quad
    \beta_k(t)
  \,\Bigr]^\top,
  \label{eq:res_feat}
\end{equation}
representing the battery-derated CPU ratio,
mean link-bandwidth ratio to active peers,
and residual battery fraction, respectively.
$\mathbf{E}^R$ is the edge index of the fully-connected resource
graph, fixed for a given fleet size $M$.

\textit{Scheduling feedback.}
$\mathbf{f}_t \in \mathbb{R}^6$ summarises the outcome of the
most recent scheduling round:
\begin{equation}
  \mathbf{f}_t = \bigl[
    \lambda_t,\;\sigma_t,\;\varrho_t,\;
    \bar{u}_t,\;\hat{B}_{\min,t},\;\hat{\beta}_{\min,t}
  \bigr]^\top,
  \label{eq:feedback}
\end{equation}
where $\lambda_t = T_{\mathrm{span}}(t)/T_D$ is the normalised
makespan, $\sigma_t$ is the task-success rate, $\varrho_t$ is the
number of forced reschedule events, $\bar{u}_t$ is the mean fleet
CPU utilisation, $\hat{B}_{\min,t}$ is the normalised minimum active
link bandwidth, and $\hat{\beta}_{\min,t}$ is the minimum battery
fraction across all UAVs.
Together, these six signals equip the agent with a compressed
view of both the scheduling quality and the current physical
bottleneck of the swarm.

\subsubsection*{Action Space}

The action space is a flat discrete set of
\emph{node-level graph-shaping operations}:
\begin{equation}
  \mathcal{A} = \mathcal{A}_{\mathrm{noop}}
              \;\cup\; \mathcal{A}_{\mathrm{split}}
              \;\cup\; \mathcal{A}_{\mathrm{merge}},
  \label{eq:action_space}
\end{equation}
with cardinalities
$|\mathcal{A}_{\mathrm{noop}}| = 1$,\,
$|\mathcal{A}_{\mathrm{split}}| = N_{\max}$, and
$|\mathcal{A}_{\mathrm{merge}}| = \binom{N_{\max}}{2}$,
giving $|\mathcal{A}| = 1 + 20 + 190 = 211$ for $N_{\max} = 20$.
Each action is decoded to a triple
$(op,\, i,\, j) \in
 \{\mathrm{noop},\mathrm{split},\mathrm{merge}\}
 \times [N_{\max}] \times [N_{\max}]$,
where $i$ and $j$ index nodes in the padded ordering.
Operating at the node level, rather than selecting a global operation
heuristically, allows the policy to learn \emph{which specific task}
to restructure based on the current resource context—a capability
absent from prior graph-shaping approaches.

\subsubsection*{Reward Function}

The reward $r_t$ is a composite signal designed to credit
absolute scheduling quality, incremental improvement, and
penalise structural pathologies:
\begin{equation}
  r_t = r_t^{\mathrm{qual}}
      + r_t^{\Delta}
      + r_t^{\mathrm{pen}},
  \label{eq:reward}
\end{equation}
with each component defined as follows.

The \textit{absolute-quality term} rewards the agent for achieving
high success rates and feasible latencies, irrespective of prior
performance:
\begin{equation}
  r_t^{\mathrm{qual}} =
    2\,\sigma_t + 1.5\,\max\!\bigl(0,\,1-\lambda_t\bigr).
  \label{eq:r_qual}
\end{equation}

The \textit{improvement term} provides dense shaping by rewarding
reductions in normalised latency and increases in success rate
relative to the preceding epoch:
\begin{equation}
  r_t^{\Delta} =
    (\lambda_{t-1} - \lambda_t)
    + 0.5\,(\sigma_t - \sigma_{t-1}).
  \label{eq:r_delta}
\end{equation}

The \textit{penalty term} discourages forced reschedules,
deadline violations, and unbounded graph inflation:
\begin{equation}
  r_t^{\mathrm{pen}} =
    -\,\frac{\varrho_t}{N_t}
    - 2\,\max\!\bigl(0,\,\lambda_t - 1\bigr)
    - 0.05\,\max\!\bigl(0,\,N_t - (N_{\max}-2)\bigr).
  \label{eq:r_pen}
\end{equation}
The reschedule penalty is normalised by the current graph size $N_t$
to eliminate the bias whereby larger graphs accumulate higher absolute
reschedule counts by construction.
The latency overshoot penalty is smooth rather than clipped,
preserving gradient signal beyond the deadline boundary.
The inflation penalty activates only when $N_t > 18$, preserving
headroom for productive splitting while preventing degenerate
graph growth.

% ----------------------------------------------------------------
\subsection{Dual-Graph State Representation}
\label{subsec:gnn}

A dedicated GNN encoder extracts structured representations from
each graph independently before the two streams are fused.

\subsubsection*{Node-Level GraphSAGE Encoding}

Both graphs are processed by an instance of the same
\textit{node-level} GNN architecture, which differs from standard
graph classification networks by omitting the global pooling step
and returning a per-node embedding matrix.
Concretely, for a graph with node features $\mathbf{X}$ and edge
index $\mathbf{E}$, the encoder applies $L$ GraphSAGE
layers~\cite{hamilton2017inductive} followed by a linear projection:
\begin{equation}
  \mathbf{h}_i^{(\ell)} = \sigma\!\Bigl(
    \mathbf{W}^{(\ell)} \cdot \mathrm{CONCAT}\Bigl[
      \mathbf{h}_i^{(\ell-1)},\;
      \mathrm{MEAN}_{j \in \mathcal{N}(i)}\!\bigl[
        \mathbf{h}_j^{(\ell-1)}
      \bigr]
    \Bigr]
  \Bigr),
  \label{eq:sage}
\end{equation}
where $\mathcal{N}(i)$ is the neighbourhood of node $i$
and $\sigma$ denotes LayerNorm followed by ReLU.
For the task graph and resource graph respectively:
\begin{align}
  \mathbf{H}^T &= \textsc{Enc}_\tau\!\bigl(
                   \mathbf{X}^T_t,\,\mathbf{E}^T_t
                 \bigr) \in \mathbb{R}^{N_{\max} \times d},
  \label{eq:enc_task}\\
  \mathbf{H}^R &= \textsc{Enc}_r\!\bigl(
                   \mathbf{X}^R_t,\,\mathbf{E}^R
                 \bigr) \in \mathbb{R}^{M \times d},
  \label{eq:enc_res}
\end{align}
with embedding dimension $d = 32$, hidden width $d_h = 64$,
and $L = 2$ message-passing layers.
Padding rows in $\mathbf{H}^T$ are zeroed out after encoding via a
boolean mask derived from $\mathbf{X}^T_t$, preventing spurious
gradient flow through phantom nodes.

% ----------------------------------------------------------------
\subsection{Cross-Graph Attention Mechanism}
\label{subsec:attention}

The central architectural innovation of RTGS is a
\emph{cross-graph attention} module that enriches each task node
embedding with contextual information drawn from the current UAV
resource state.
This coupling enables the policy to condition split and merge
decisions on \emph{which physical resources would be responsible for
executing the affected tasks}, a form of lookahead that is
inaccessible to policies that process the two graphs independently.

\subsubsection*{Task-to-Resource Attention}

Given the encoded representations $\mathbf{H}^T$ and $\mathbf{H}^R$,
single-head scaled dot-product attention is applied with task nodes
as queries and resource nodes as keys and values:
\begin{align}
  \mathbf{Q} &= \mathbf{H}^T \mathbf{W}_Q,\quad
  \mathbf{K}  = \mathbf{H}^R \mathbf{W}_K,\quad
  \mathbf{V}  = \mathbf{H}^R \mathbf{W}_V,
  \label{eq:qkv}\\[4pt]
  \boldsymbol{\Alpha} &=
    \mathrm{softmax}\!\left(
      \frac{\mathbf{Q}\mathbf{K}^\top}{\sqrt{d}}
    \right) \in [0,1]^{N_{\max} \times M},
  \label{eq:attn}\\[4pt]
  \tilde{\mathbf{H}}^T &=
    \mathbf{H}^T + \mathbf{W}_O \bigl(\boldsymbol{\Alpha}\mathbf{V}\bigr),
  \label{eq:enriched}
\end{align}
where $\mathbf{W}_Q, \mathbf{W}_K, \mathbf{W}_V, \mathbf{W}_O
\in \mathbb{R}^{d \times d}$ are learnable projection matrices
and the outer sum in \eqref{eq:enriched} is a residual connection.

The entry $\Alpha_{ik} \in [0,1]$ can be interpreted as the
\emph{soft assignment weight} of task $v_i$ to UAV $u_k$:
a high value indicates that the current resource embedding of $u_k$
is most compatible with executing $v_i$.
These weights are not discarded after computing $\tilde{\mathbf{H}}^T$;
they are propagated to the merge scoring head as an inter-node
link-quality signal (see Section~\ref{subsec:actor}).

\subsubsection*{Global State Vector}

A global state vector $\mathbf{g} \in \mathbb{R}^{d_g}$
is assembled for use by the value function and the noop scoring head:
\begin{equation}
  \mathbf{g} = \Bigl[
    \bar{\mathbf{h}}^T_{\mathrm{real}}\;\Big\|\;
    \bar{\mathbf{h}}^R\;\Big\|\;
    \mathbf{b}_{\mathrm{res}}\;\Big\|\;
    \mathbf{f}_t
  \Bigr],
  \label{eq:global}
\end{equation}
where $\bar{\mathbf{h}}^T_{\mathrm{real}}$ is the masked mean of
$\tilde{\mathbf{H}}^T$ over real (non-padding) task nodes,
$\bar{\mathbf{h}}^R$ is the mean of $\mathbf{H}^R$ over all UAVs,
and $\mathbf{b}_{\mathrm{res}} \in \mathbb{R}^6$ is a
\textit{resource bottleneck descriptor}:
\begin{equation}
  \mathbf{b}_{\mathrm{res}} =
  \Bigl[
    \min_{k}\,\mathbf{x}_k^R(t)\;\Big\|\;
    \tfrac{1}{M}\sum_k \mathbf{x}_k^R(t)
  \Bigr],
  \label{eq:bottleneck}
\end{equation}
concatenating the per-feature fleet minimum and fleet mean across
all $M$ UAVs.
The minimum captures the identity of the most depleted or
most congested node in the fleet—information that is lost under
plain mean aggregation—so that the value function remains
sensitive to impending resource failures.
With $d = 32$, $M = 8$, and a six-dimensional feedback vector,
the global state has dimension
$d_g = d + d + 6 + 6 = 76$.

% ----------------------------------------------------------------
\subsection{Node-Level Actor-Critic}
\label{subsec:actor}

\subsubsection*{Three-Head Action Logit Architecture}

The policy network produces action logits through three
semantically distinct scoring heads, each operating at a
different level of granularity.

\noindent\textbf{Noop head.}
A linear map over the global state $\mathbf{g}$ produces a single
scalar logit representing the value of leaving the graph unchanged:
\begin{equation}
  z_{\mathrm{noop}} = \mathbf{w}_{\mathrm{no}}^\top \mathbf{g}
    \in \mathbb{R}.
  \label{eq:z_noop}
\end{equation}

\noindent\textbf{Split head.}
A shared linear scoring function is applied independently to each
resource-enriched task node embedding, yielding one logit per
padded node position:
\begin{equation}
  \mathbf{z}_{\mathrm{split}} =
    \tilde{\mathbf{H}}^T\,\mathbf{w}_{\mathrm{sp}}
    \in \mathbb{R}^{N_{\max}}.
  \label{eq:z_split}
\end{equation}
Because $\tilde{\mathbf{H}}^T$ already encodes the soft assignment
of each task to a UAV, the split head learns to prefer nodes whose
assigned UAV has high residual CPU capacity—a form of
\emph{resource-opportunistic decomposition}.

\noindent\textbf{Merge head with link-quality signal.}
For each ordered pair $(v_i, v_j)$ in the enumeration of all
$\binom{N_{\max}}{2} = 190$ candidate node pairs, a
\textit{link-quality scalar} $q_{ij}$ is computed from the
cross-graph attention weights:
\begin{equation}
  q_{ij} = \boldsymbol{\Alpha}_i \cdot \boldsymbol{\Alpha}_j
          = \sum_{k=1}^{M} \Alpha_{ik}\,\Alpha_{jk},
  \label{eq:link_quality}
\end{equation}
where $\boldsymbol{\Alpha}_i \in \mathbb{R}^M$ is the $i$-th row of
the attention matrix.
A high value of $q_{ij}$ indicates that both tasks attend strongly to
the \emph{same} UAV, implying that executing them on the same node
would incur negligible transfer cost—and that merging them
provides limited scheduling benefit.
Conversely, a low $q_{ij}$ signals that the two tasks are likely
assigned to \emph{different} UAVs connected by a potentially weak
link; merging eliminates that cross-UAV transfer entirely and
is therefore highly advantageous.
The merge logit is computed by a two-layer MLP that jointly
processes both node embeddings and the link-quality signal:
\begin{equation}
  z_{ij}^{\mathrm{merge}} =
    \mathrm{MLP}_{\mathrm{me}}\!\Bigl(
      \bigl[\,
        \tilde{\mathbf{h}}_i \;\|\;
        \tilde{\mathbf{h}}_j \;\|\;
        q_{ij}
      \,\bigr]
    \Bigr) \in \mathbb{R},
  \label{eq:z_merge}
\end{equation}
with input dimension $2d + 1 = 65$.

All three logit vectors are concatenated into the full action score:
\begin{equation}
  \mathbf{z} = \bigl[
    z_{\mathrm{noop}}\;\|\;
    \mathbf{z}_{\mathrm{split}}\;\|\;
    \mathbf{z}_{\mathrm{merge}}
  \bigr] \in \mathbb{R}^{211}.
  \label{eq:logits}
\end{equation}
The value estimate is obtained from a separate linear head over
the global state:
\begin{equation}
  V(s_t) = \mathbf{w}_V^\top \mathbf{g} \in \mathbb{R}.
  \label{eq:value}
\end{equation}

\subsubsection*{Structural Action Masking}
\label{subsubsec:masking}

A key property of RTGS is that the action distribution is restricted
to \emph{structurally valid} operations at every step by applying a
Boolean mask $\mathbf{m}_t \in \{0,1\}^{211}$ before sampling.
The mask is computed from the current observation without gradient
and enforces three validity conditions:

\begin{enumerate}[label=(\roman*), leftmargin=*, itemsep=2pt]
  \item \textbf{Noop} is always valid:
        $m_0 = 1$.
  \item \textbf{Split}$(v_i)$ is valid iff $v_i$ is a real
        (non-padding) node, the graph has not reached its maximum
        cardinality ($N_t < N_{\max}$), and at least one UAV
        retains positive effective CPU capacity
        ($\max_k \tilde{\phi}_k(t) > 0$):
        \begin{equation}
          m_i^{\mathrm{sp}} =
            \mathbb{1}[v_i\text{ real}] \cdot
            \mathbb{1}[N_t < N_{\max}] \cdot
            \mathbb{1}\!\bigl[\max_k \tilde{\phi}_k(t) > 0\bigr].
          \label{eq:mask_split}
        \end{equation}
  \item \textbf{Merge}$(v_i, v_j)$ is valid iff both nodes
        are real and are \emph{leaf} nodes in the current DAG
        (out-degree zero):
        \begin{equation}
          m_{ij}^{\mathrm{me}} =
            \mathbb{1}[v_i\text{ real leaf}] \cdot
            \mathbb{1}[v_j\text{ real leaf}].
          \label{eq:mask_merge}
        \end{equation}
\end{enumerate}

Condition (iii) guarantees DAG acyclicity: if both $v_i$ and $v_j$
are leaves, no directed path exists between them, so their merger
cannot introduce a cycle.
The masked Categorical distribution is formed as:
\begin{equation}
  \pi_\theta(a \mid s_t) =
    \frac{\exp(z_a + \log m_{t,a})}
         {\sum_{a'} \exp(z_{a'} + \log m_{t,a'})},
  \label{eq:masked_policy}
\end{equation}
where $\log 0 \triangleq -\infty$, so that invalid actions
receive zero probability.
The mask is reapplied identically during PPO updates to maintain
consistency between rollout-collection and gradient-computation
log-probabilities.

% ----------------------------------------------------------------
\subsection{Graph Shaping Operators}
\label{subsec:operators}

Two graph-transformation primitives are executed by the environment
in response to the agent's action.
Both operators produce a mutated copy of $\mathcal{G}_T$ without
modifying the original, so that any roll-back or comparison is
well-defined.

\subsubsection*{Split Operator}

\textsc{Split}$(v_i)$ decomposes a targeted task node into
two parallel sub-tasks.
Formally, given $|\mathcal{V}_T| < N_{\max}$, a sibling node
$v_i^{+}$ is introduced with
\begin{equation}
  c_{i^{+}} = c_i/2,\qquad s_{i^{+}} = s_i/2,
  \label{eq:split_attr}
\end{equation}
and the attributes of $v_i$ are simultaneously halved.
The sibling inherits the complete in-neighbourhood and
out-neighbourhood of $v_i$, while no edge is inserted between
$v_i$ and $v_i^{+}$:
\begin{equation}
  \mathcal{E}_T' = \mathcal{E}_T
    \cup \bigl\{(v_p,\, v_i^{+}) : v_p \in \Pi(v_i)\bigr\}
    \cup \bigl\{(v_i^{+},\, v_s) : v_s \in \Sigma(v_i)\bigr\}.
  \label{eq:split_edges}
\end{equation}
The absence of a sibling edge is intentional:
$v_i$ and $v_i^{+}$ are topologically independent and may therefore
execute concurrently on separate UAVs,
halving the effective CPU bottleneck of the original task and
exposing additional inter-UAV parallelism.

\subsubsection*{Merge Operator}

\textsc{Merge}$(v_i, v_j)$ collapses two tasks into a single node.
The acyclicity guard
\begin{equation}
  \nexists\;\text{directed path}\;v_i \leadsto v_j
  \;\;\text{and}\;\;
  \nexists\;\text{directed path}\;v_j \leadsto v_i
  \label{eq:cycle_guard}
\end{equation}
is verified before the operation; otherwise the action is silently
voided.
When the guard holds, the attributes of $v_j$ are absorbed into
$v_i$:
\begin{equation}
  c_i \leftarrow c_i + c_j,\qquad s_i \leftarrow s_i + s_j.
  \label{eq:merge_attr}
\end{equation}
All edges formerly incident to $v_j$ are redirected to $v_i$
without duplicating existing edges or introducing self-loops;
$v_j$ is then removed.
The critical scheduling benefit is the elimination of the
inter-UAV data transfer $d_{ij}$ that would have been incurred on
the link connecting the respective executors of $v_i$ and $v_j$—a
saving that is particularly significant when that link suffers from
mobility-induced bandwidth degradation.

% ----------------------------------------------------------------
\subsection{PPO-Based Policy Optimisation}
\label{subsec:ppo}

The policy $\pi_\theta$ is trained with Proximal Policy
Optimisation (PPO)~\cite{schulman2017proximal} using a rollout
buffer of $T_r = 2048$ transitions collected from a single
environment instance.
The compound loss is
\begin{equation}
  \mathcal{L}(\theta) =
    \mathcal{L}_{\mathrm{clip}}(\theta)
    + c_v \mathcal{L}_{\mathrm{vf}}(\theta)
    - c_H \mathcal{H}(\pi_\theta),
  \label{eq:ppo_loss}
\end{equation}
where the clipped policy gradient objective is
\begin{equation}
  \mathcal{L}_{\mathrm{clip}}(\theta) =
    -\mathbb{E}_t\!\left[
      \min\!\Bigl(
        \rho_t \hat{A}_t,\;
        \mathrm{clip}\bigl(\rho_t,\,1\!-\!\epsilon,\,1\!+\!\epsilon\bigr)\hat{A}_t
      \Bigr)
    \right],
  \label{eq:ppo_clip}
\end{equation}
with likelihood ratio
$\rho_t = \pi_\theta(a_t \mid s_t)\,/\,\pi_{\theta_{\mathrm{old}}}(a_t \mid s_t)$
and clip coefficient $\epsilon = 0.2$.
The value loss and entropy regulariser are
\begin{align}
  \mathcal{L}_{\mathrm{vf}}(\theta)
    &= \mathbb{E}_t\!\bigl[\bigl(V_\theta(s_t) - R_t\bigr)^2\bigr],
  \label{eq:vf_loss}\\
  \mathcal{H}(\pi_\theta)
    &= -\mathbb{E}_t\!\bigl[
        \sum_a \pi_\theta(a \mid s_t)\log\pi_\theta(a \mid s_t)
      \bigr],
  \label{eq:entropy}
\end{align}
with discount-return $R_t = \sum_{l \geq 0} \gamma^l r_{t+l}$,
$\gamma = 0.99$, and coefficients $c_v = 0.5$, $c_H = 0.05$.
The entropy coefficient is set higher than typical PPO defaults to
counteract the tendency of the policy to prematurely collapse onto
the noop action in the large 211-dimensional action space.
The advantage estimates $\hat{A}_t = R_t - V_\theta(s_t)$ are
standardised to zero mean and unit variance per rollout.
Parameters are updated for $K = 4$ epochs per rollout with
mini-batches of size 128; gradients are clipped to $\ell_2$
norm~$0.5$.

The full training procedure is summarised in
Algorithm~\ref{alg:rtgs}.

% ----- Algorithm block -----------------------------------------
\begin{algorithm}[t]
\caption{RTGS Training with PPO}
\label{alg:rtgs}
\begin{algorithmic}[1]
\Require Environment $\mathcal{E}$, policy network $\pi_\theta$
         with parameters $\theta$,\;
         rollout length $T_r$,\;
         PPO epochs $K$,\;
         mini-batch size $B_{\mathrm{mb}}$
\State Initialise $\theta$ randomly; reset $\mathcal{E}$; receive $s_1$

\While{total steps $< 2{\times}10^6$}

  \State \textbf{// --- Phase 1: Rollout collection ---}
  \State Clear rollout buffer $\mathcal{B}$
  \For{$t = 1, \ldots, T_r$}

    \State Encode observation:
           $\mathbf{H}^T, \mathbf{H}^R \leftarrow
            \textsc{Enc}_\tau(s_t),\,\textsc{Enc}_r(s_t)$
    \State Compute cross-attention:
           $\tilde{\mathbf{H}}^T,\, \boldsymbol{\Alpha}
            \leftarrow \textsc{CrossAttn}(\mathbf{H}^T, \mathbf{H}^R)$
    \State Build global state $\mathbf{g}$ via \eqref{eq:global}
    \State Compute logits $\mathbf{z}$ via \eqref{eq:logits}
           and value $V(s_t)$ via \eqref{eq:value}
    \State Compute action mask $\mathbf{m}_t$ via
           \eqref{eq:mask_split}--\eqref{eq:mask_merge}
    \State Sample $a_t \sim \mathrm{Categorical}(
           \mathrm{softmax}(\mathbf{z} + \log\mathbf{m}_t))$
    \State Decode $a_t \rightarrow (op,\, v_i,\, v_j)$
    \State Execute in $\mathcal{E}$:
           $s_{t+1},\, r_t,\, \mathrm{done}_t
            \leftarrow \mathcal{E}.\textsc{Step}(a_t)$
           \Comment{shaping $\rightarrow$ scheduling $\rightarrow$ mobility}
    \State Store $(s_t, a_t, \log\pi_\theta(a_t|s_t),
                  r_t, V(s_t), \mathrm{done}_t)$ in $\mathcal{B}$

  \EndFor

  \State \textbf{// --- Phase 2: Return and advantage estimation ---}
  \State $R_{T_r+1} \leftarrow V_\theta(s_{T_r+1})$
         (bootstrap from last state)
  \State Compute discounted returns $R_t$ and standardised advantages
         $\hat{A}_t$ for all $t$

  \State \textbf{// --- Phase 3: Policy update ---}
  \State Snapshot old policy: $\theta_{\mathrm{old}} \leftarrow \theta$
  \For{epoch $= 1, \ldots, K$}
    \For{each mini-batch $\mathcal{M} \subseteq \mathcal{B}$
         of size $B_{\mathrm{mb}}$}
      \State Recompute $\mathbf{z}, V_\theta(s_t)$ with current $\theta$
      \State Reapply action mask $\mathbf{m}_t$; compute
             $\log\pi_\theta(a_t|s_t)$
      \State Compute $\mathcal{L}(\theta)$ via
             \eqref{eq:ppo_loss}--\eqref{eq:entropy}
      \State Update: $\theta \leftarrow \theta - \eta
             \,\nabla_\theta\mathcal{L}(\theta)$
             with gradient clipping ($\|\nabla\|_2 \leq 0.5$)
    \EndFor
  \EndFor

\EndWhile
\State \Return trained policy $\pi_\theta$
\end{algorithmic}
\end{algorithm}

\textbf{Remark on the environment step.}
The call $\mathcal{E}.\textsc{Step}(a_t)$ encapsulates a
six-stage pipeline:
(i) decode and apply the shaping operator to $\mathcal{G}_T$;
(ii) run the mobility-aware greedy scheduler on the modified
graph $\mathcal{G}_T'$ with current $\mathbf{B}(t)$
and $\tilde{\boldsymbol{\phi}}(t)$;
(iii) deduct battery according to \eqref{eq:energy_depletion};
(iv) advance UAV positions via the Random Waypoint model;
(v) compute the reward $r_t$;
(vi) construct the updated observation $s_{t+1}$.
This deterministic, non-differentiable environment step is the
reason RL is necessary: the combinatorial scheduling layer prevents
end-to-end gradient propagation from the task-assignment outcome
back through the graph-shaping decision.
